% =====================================================================
%  23.5  The Engine Boundary: choosing among Li Chao / beats / ODT /
%        composition -- and when an EARLIER chapter's engine is the right one.
%  WORKED EXAMPLES ONLY -- NO \solvedproblem (in-section count stays 4).
%  Exports: sec:c29-boundary  def:c29-map  lst:c29-map
%  DISCHARGES (prose + \chref/\secref only, NO re-link):
%    * Ch.20 \secref{c15-strains} (20.5): "chmin/chmax/modulo is segment-tree
%      beats, Ch.23, no code there" -> delivered in \secref{c29-beats}.
%    * Ch.21 \secref{c19-sparse} (21.5): forward-cited advanced trees -> Li Chao
%      (\secref{c29-lichao}) and beats live here.
%    * Ch.22 \secref{c16-boundary} case (C): "a path chmin needs a beats node,
%      not Ch.20's tag (Ch.23)" and its lastword "Next: c29" -> this chapter.
%  Stays OUT (name + forward-cite, no code): fully DYNAMIC convex hull / Li Chao
%    WITH DELETIONS -> link-cut / Euler-tour Ch.25 (\chref{c23}); offline
%    reorderable queries -> Mo's Ch.27 (\chref{c21}).
%  Do NOT \term anything already indexed elsewhere: "segment-tree beats" (Ch.20
%    \secref{c15-strains} AND Ch.21 \secref{c19-sparse}), "Li Chao" (Ch.21), the
%    lazy/composition/monoid words (Ch.20); and this chapter's own Li Chao line
%    container / Chtholly tree / function composition (00/01/03/04).  The
%    Ch.22 INDEX-overfull lesson: a redundant \term overflows the doc-wide INDEX.
% =====================================================================
\section{The Engine Boundary: Which Advanced Tree to Reach For}\label{sec:c29-boundary}

\begin{previouslybox}[Four sections taught three engines; this one says when NOT to use them]
Sections 23.1--23.4 built the Li Chao line container, its tree-riding variants,
segment-tree beats and the Chtholly tree, and the function-composition tree.  Each
rescues a family of operations that an earlier chapter's engine deliberately
\emph{refused}: Chapter~19 (\(\secref{c13-static}\)) folds a fixed array,
Chapter~20 (\(\secref{c15-strains}\)) defers only \emph{composable} tags,
Chapter~21 (\(\secref{c19-sparse}\)) optimises the fully static case, and
Chapter~22 (\(\secref{c16-boundary}\)) reduces a tree path to array ranges but
insists the per-range op be a monoid.  Every one of them forward-cited \chref{c29}
for a non-composable or envelope-shaped operation.  Those promises are discharged
below.  A boundary section earns its place because these four engines are
temptingly general -- a solver who has just learned beats will try to beat every
range update -- and each has a precise cost and a failure mode that a plainer
engine avoids.
\end{previouslybox}

\begin{topicintro}
Four diagnostic questions place any range/array-with-DP problem on this chapter's
map.  \emph{What does a node aggregate?} A single number (sum/min/gcd) is
Chapter~19's fold; a \emph{line envelope} is Li Chao; a \emph{clamp bookkeeping}
triple is beats; a \emph{transform} is a composition tree.  \emph{What is the
domain?} Array positions favour every engine; a numeric \(x\)-domain (\(\min/\max\)
of lines at a queried \(x\)) forces Li Chao.  \emph{Do the updates compose into one
pending value?} If yes, Chapter~20's lazy tag is strictly cheaper than beats --
beats is only worth its amortised machinery when they do not.  \emph{Worst case or
random?} The same clamp is beats (worst-case \(O(\log n)\)) or a Chtholly tree
(expected-case, twenty lines) depending on whether you trust the test data.

\textbf{Prerequisites.} The whole of 23.1--23.4; Ch.~19 (\(\chref{c13}\)) for the
fold these all extend; Ch.~20 (\(\secref{c15-compose}\)) for the composition monoid
that a lazy tag and a composition tree share; Ch.~22 (\(\secref{c16-boundary}\))
for the path-to-range reduction that can \emph{wrap} any engine here.
\end{topicintro}

% ---------------------------------------------------------------
\subsection{The axis that decides the engine: what a node stores}
% ---------------------------------------------------------------

\begin{definition}{The advanced-array engine map}{c29-map}
Classify by \emph{what the node's aggregate is} and \emph{what the update does to
it}.  \(\{\text{number}, \text{composable update}\}\) is Chapter~20's lazy tree.
\(\{\text{number}, \text{data-dependent clamp/modulo}\}\) is segment-tree beats.
\(\{\text{line envelope over an }x\text{-domain}, \text{insert/query}\}\) is Li
Chao.  \(\{\text{transform}, \text{compose on merge}\}\) is a function-composition
tree.  \(\{\text{equal-runs}, \text{random data}\}\) is a Chtholly tree.  These
five cells plus the plain fold and heavy-light are the whole family; each row of
the table below names the tell that selects one.
\end{definition}

\begin{center}
\small
\renewcommand{\arraystretch}{1.26}
\setlength{\tabcolsep}{4pt}
\begin{tabular}{@{}p{0.28\linewidth} p{0.24\linewidth} >{\raggedright\arraybackslash}p{0.20\linewidth} >{\centering\arraybackslash}p{0.18\linewidth}@{}}
\toprule
\textbf{Symptom in the spec} & \textbf{Right engine} & \textbf{Why the others lose} & \textbf{Where}\\
\midrule
fold a number over a range; point updates & plain segment tree & any extra state is dead weight & Ch.~19 \(\chref{c13}\)\\
range add/assign/affine + range query & lazy tag & composes into one pending value -- beats' descend is waste & Ch.~20 \(\chref{c15}\)\\
insert lines, query best at point \(x\); arbitrary slopes & Li Chao container & a position-tree has no \(x\)-axis to fold & \(\secref{c29-lichao}\)\\
DP over ancestors / descendants of a tree, transitions are lines & Li Chao on tree (rollback / merge) & heavy-light folds values, not envelopes & \(\secref{c29-path}\)\\
range \(\mathit{chmin}/\mathit{chmax}\)/modulo \emph{and} sum, worst case & segment-tree beats & a lazy tag cannot summarise a clamp & \(\secref{c29-beats}\)\\
same clamp ops but data/queries are RNG-generated & Chtholly tree (ODT) & beats is more code; ODT is short and expected-fast & \(\defref{c29-odt}\)\\
range = run a fixed-state automaton / recurrence & composition tree (min--plus matrix) & a scalar fold cannot hold a state machine & \(\secref{c29-func}\)\\
path update/query on a tree, op composable & heavy-light + one engine here & (HLD reduces to ranges; the leaf engine is still chosen above) & Ch.~22 \(\chref{c16}\)\\
\bottomrule
\end{tabular}
\end{center}

\begin{lstlisting}[language={}, basicstyle=\ttfamily\small, numbers=none, frame=single,
  caption={Decision ladder for the advanced-array family. Read top to bottom, stop at the first line that matches. Hands off to Chapters 19/20/21/22 whenever a plainer engine already owns the cell.}, label={lst:c29-map}]
  Is the query just "fold a stored number over a range" (sum/min/gcd/count)?
    |
    +-- YES, updates are single positions -> plain SEGMENT TREE (Ch.19, c13).
    +-- YES, but whole ranges get ADD / ASSIGN / AFFINE, online:
    |         -> LAZY PROPAGATION (Ch.20, c15).  It composes into one tag.
    |            (Do NOT reach for beats here -- the extra state is pure cost.)
    |
    +-- What UPDATE shape breaks the "one tag" rule?
    |
    |   (a) a data-DEPENDENT clamp: range chmin/chmax/modulo alongside a sum/max.
    |         -> SEGMENT-TREE BEATS (this ch., 23.3).  Store (max,smax,cnt),
    |            prune / settle / descend; amortised O(log n) worst case.
    |         -> same but the input is RANDOM (generated ops, few adversarial)?
    |            a CHTHOLLY TREE (23.3) is 20 lines and expected-fast -- but it
    |            TLEs if the clamp-free tests appear (see 896C's documented hack).
    |
    |   (b) the aggregate is an ENVELOPE of functions, queried at a chosen x:
    |         "insert a line, ask min/max at x"  -> LI CHAO (23.1).
    |         that query, over a tree's ancestors -> Li Chao + ROLLBACK (23.2);
    |         over a subtree -> Li Chao MERGE small-to-large (23.2), e.g. 932F.
    |
    |   (c) the range query runs a SEQUENTIAL machine (automaton / recurrence)
    |       with a FIXED small state  -> COMPOSITION TREE of (min,+) matrices
    |       (23.4): store the composed map, merge by matrix product. 750E.
    |
    +-- Is it on a TREE PATH?  Reduce with HEAVY-LIGHT (Ch.22, c16) to O(log n)
         ranges, THEN plug in whichever engine (a)-(c) the row above selects.

  STILL OUT OF SCOPE here (forward-cites only, no code):
    * hull that DELETES / arbitrary insert order -> DYNAMIC convex hull via
      link-cut / Euler-tour trees (Ch.25, c23).
    * queries you may REORDER, aggregate has no inverse -> Mo's on arrays/trees
      (Ch.27, c21).
    * topology or root CHANGES between edits -> link-cut (Ch.25, c23), not any
      static decomposition.
\end{lstlisting}

% ---------------------------------------------------------------
\subsection{Worked boundary calls (read a spec, name the engine)}
% ---------------------------------------------------------------

\begin{comparisonbox}[Four specifications, four different answers -- none solved here]
\textbf{(A) ``\(n,q\le10^5\); operation `set \(a_i\leftarrow\min(a_i,x)\) on a range', `add \(x\) on a range'; query the range sum worst-case.''}
A clamp is \emph{not} a monoid action -- Chapter~20's \(\secref{c15-strains}\) proved
the fold is unrecoverable from \((\mathit{sum},\max)\) alone, which is precisely the
forward-cite that lands on this chapter.  Segment-tree beats
(\(\secref{c29-beats}\)) with the \((\max,\mathrm{smax},\mathrm{cnt})\) node is the
answer; the add composes as an ordinary lazy add on top.  Expected data would admit
a Chtholly tree, but ``worst-case'' rules it out.

\textbf{(B) ``convex-hull trick DP, \(f_i=\min_{j<i}(m_j x_i+b_j)+c_i\), slopes \(m_j\) arrive with no monotonicity.''}
The classic monotone-slope deque needs slopes and queries both sorted; arbitrary
\(m_j\) kills it.  This is the Li Chao container (\(\secref{c29-lichao}\)): index the
\(x\)-domain, insert line \(j\) as \(j\) is computed, query at \(x_i\) --
\(O(\log C)\) each.  If the same envelope must be read over a tree's ancestors
rather than a prefix, it becomes the rollback container of \(\secref{c29-path}\).

\textbf{(C) ``\(n,q\le2\cdot10^5\) digits; many queries `fewest deletions so this substring contains 2017 but not 2016'.''}
Neither a fold nor a clamp -- the answer is the cost of running a fixed
\emph{automaton} over the range.  That is a function-composition tree of
\(5\times5\) min-plus matrices (\(\secref{c29-func}\), solved as CF~750E); beats and
Li Chao have nothing to say here.

\textbf{(D) ``fixed weighted tree; `lower every node value on path \(u\)--\(v\) toward \(x\)'; query path sum.''}
Chapter~22's \(\secref{c16-boundary}\) case~(C) deferred exactly this: keep the
heavy-light chain-range reduction, but replace the leaf engine with a beats node.
So the answer is HLD \emph{wrapping} \(\secref{c29-beats}\) -- the two engines
composed, not competing.  This discharges that forward-cite and the closing
``Next: \chref{c29}'' of Chapter~22.
\end{comparisonbox}

\begin{observation}[The engines are chosen by the \emph{aggregate}, orthogonal to the \emph{container}]
Li Chao and beats both build on the Chapter~19 skeleton; they differ only in what
a node keeps (a line vs.\ a max-triple) and what ``combine'' means (pointwise best
vs.\ pushup of three values).  A composition tree keeps a \emph{map} and composes.
Hence ``tree path'' (Ch.~22) and ``which monoid'' (this chapter) are \emph{two
independent} decisions: heavy-light reduces a path to ranges, and then the range
engine is picked from \(\defref{c29-map}\) exactly as it would be on a plain array.
Confusing the axis of the container with the axis of the aggregate is the most
common mis-selection in this family.
\end{observation}

\begin{notebox}[Beats versus the Chtholly tree: the honest tie-break]
They solve the \emph{same} operations (add, \(\mathit{chmin}\), assign, order
statistics).  Beats gives a worst-case \(O((n+q)\log n)\)
(\(\thmref{c29-beats-cost}\)) and needs the \((\max,\mathrm{smax},\mathrm{cnt})\)
machinery.  The Chtholly tree (\(\defref{c29-odt}\)) is a dozen lines and expected
\(O((n+q)\log n)\) \emph{on random data}, because each \(\mathit{chmin}\) collapses
runs -- but CF~896C's history shows a single \(\mathit{chmin}\)-free adversarial test
drives the run count to \(n\) and the ODT to \(\tle\).  Default to beats for
worst-case correctness; keep the ODT for generated/random instances or as a fast
prototype you can beat later.
\end{notebox}

\begin{stuckbox}[If you reach for beats but the update actually composes]
Write the update as a function and check whether two of them collapse to one.
\(\mathit{add}\), assign, and affine \(x\mapsto ax+b\) all do -- that is a lazy tag
(Ch.~20), cheaper than beats.  Only when the effect depends on each element's
\emph{current value relative to the bound} (\(\min/\max\) with a bound, modulo,
``divide'') does the operation refuse to compose, and only then does beats earn its
amortised bound.  Diagnosing \emph{composable vs.\ data-dependent} is the single
distinction this boundary teaches; everything else is an implementation detail of
whichever side of that line the operation falls on.
\end{stuckbox}

\begin{recallbox}[23.5 in one breath]
\emph{Pick by what the node stores: number+composable update \(\to\) lazy
(Ch.~20); number+clamp/modulo \(\to\) beats, or a Chtholly tree if the data is
random (23.3); a line envelope over an \(x\)-domain \(\to\) Li Chao (23.1), on a
tree \(\to\) rollback/merge (23.2); a fixed-state automaton \(\to\) composition
tree (23.4).  Heavy-light (Ch.~22) reduces a path to ranges and then plugs in one
of these -- the container and the aggregate are separate axes.  Discharges:
Ch.~20's clamp boundary, Ch.~21's advanced-tree forward-cite, and Ch.~22's case~(C)
path-\(\mathit{chmin}\) all land here.  Dynamic hulls with deletions \(\to\)
link-cut (Ch.~25); offline reorderable queries \(\to\) Mo (Ch.~27).}
\selftest{\item A range update is \(a_i\leftarrow a_i+b\) then a query asks the range
sum.  Why is beats \emph{wrong} here? \emph{(It composes into one pending add -- a
Ch.~20 lazy tag is cheaper; beats' descend/settle buys nothing.)}
\item What single test separates the Li Chao engine from a plain segment tree?
\emph{(The aggregate is a family of lines queried at an arbitrary \(x\) -- there is
an \(x\)-domain axis, not a fixed per-position value to fold.)}}
\end{recallbox}
